\section{问题三：全截面达标时间}
\subsection{全域判据与控制位置}
题目要求药材各处含水率低于0.15~kg/kg。连续下降时用等于阈值的交叉时刻作为边界时间，定义
\begin{equation}
t_d=\inf\{t\ge0:\max_{0\le r\le R_0}C(r,t)\le0.15\}.\label{eq:firstpass}
\end{equation}
程序每个内部时间层保留全网格最大值判据，并检查相邻含水率是否由内向外非增；中心控制是计算核验结果，不是仅由矩阵性质预先推出。径向离散仍不能自动保证所有连续空间点的精度，故另进行空间加密。

\subsection{达标时间与典型剖面}
在恒定末端环境延拓下，计算得到
\begin{equation}\widehat t_d=\QthreeSeconds{}~\mathrm s=\QthreeHours{}~\mathrm h.\label{eq:q3time}\end{equation}
\begin{table}[!htbp]\centering
\caption{问题3每6 h及达标时含水率（kg/kg）}
\begin{tabular}{lrrrrr}\toprule
$t$/h & 0 cm & 0.5 cm & 1.0 cm & 1.5 cm & 2.0 cm\\\midrule
6 & 1.0157 & 0.9870 & 0.9007 & 0.7548 & 0.5336\\
12 & 0.4549 & 0.4420 & 0.4021 & 0.3290 & 0.1640\\
18 & 0.2980 & 0.2907 & 0.2680 & 0.2247 & 0.0842\\
24 & 0.2371 & 0.2321 & 0.2162 & 0.1851 & 0.0663\\
30 & 0.2053 & 0.2013 & 0.1887 & 0.1637 & 0.0597\\
36 & 0.1854 & 0.1821 & 0.1714 & 0.1501 & 0.0565\\
42 & 0.1717 & 0.1687 & 0.1594 & 0.1404 & 0.0547\\
48 & 0.1615 & 0.1588 & 0.1504 & 0.1332 & 0.0536\\
54 & 0.1535 & 0.1511 & 0.1434 & 0.1275 & 0.0528\\
57.1776 & 0.1500 & 0.1477 & 0.1402 & 0.1249 & 0.0525\\
\bottomrule\end{tabular}\end{table}

中心在达标剖面中为0.1500~kg/kg，表面为\QthreeSurfaceC{}~kg/kg。图中表面曲线较早接近低含湿环境，说明仅根据表面判断会提前停机。随着含水率下降，有效扩散率中的 $e^{-0.45/C}$ 项减小，后期下降速度放缓。
\begin{figure}[!htbp]\centering\includegraphics[width=0.94\textwidth]{q3_threshold_diagnostics.pdf}
\caption{问题三的最大含水率首达过程及典型剖面}\label{fig:q3}\end{figure}
\subsection{结果文件与适用范围}
result3.xlsx从60~s起逐分钟记录至\QthreeOutputEnd{}~s，末行覆盖数值首达估计，论文末行则为事件插值剖面。首达估计与输出时刻不同，不能把末行时间直接当成更精确的达标时间。主答案使用一维中截面模型，长期二维全体积结果及其偏差另在检验节给出。
